\section*{Exclusion of Large Language Model Approaches}
\label{app:llm_exclusion}

Large language models have become an active direction in 
SVG generation. Methods like StarVector, OmniSVG, LLM4SVG, 
and InternSVG treat SVG as structured code and generate it 
through transformer architectures pretrained on large-scale 
text and image collections. Commercial services followed the 
same path: SVGMaker.io generates SVG graphics from text 
prompts using GPT, Qwen, Gemini, and other models without 
exposing the underlying architecture or training details.

SVG is sequential XML, so language models are a natural fit 
for the format. The visual quality of recent LLM-based 
outputs has improved a lot compared to earlier sequence 
models. But these approaches are excluded from the main 
analysis for one specific reason: this thesis asks how 
architectural decisions in representation and supervision 
shape the structure of generated SVGs. In optimization-based 
and dataset-driven methods, that relationship can be traced 
directly. A parametric representation produces a certain 
kind of geometry. A diffusion-based loss pushes optimization 
in a specific direction. The connection between design 
choice and structural outcome is visible.

In LLM-based generation, that connection is distributed 
across billions of parameters. The model generates SVG 
tokens, but why it chose a particular path structure or 
how it decides on primitive count cannot be isolated the 
same way. This does not make LLM approaches less valuable. 
It makes them less suited to the kind of analysis this 
thesis pursues.

Figure~\ref{fig:svgmaker} shows what this looks like in 
practice. An SVG generated by SVGMaker.io from the prompt 
``a penguin on ice with 16 shapes'' produced 24 paths in 
about 10 seconds. The rendered output looks clean. But 
open the file in a design tool and the structure tells a 
different story: several paths do not match any recognizable 
visual component, and the requested 16 shapes became 24. 
The output is still more editable than what some of the 
optimization-based models in this thesis produce, which 
says something about how fast LLM-based generation is 
improving. But the structural decisions behind it cannot 
be traced to specific architectural choices, and that is 
exactly what this thesis needs to do.

\begin{figure}[ht]
    \centering
    \begin{tabular}{cc}
    
\includegraphics[width=0.45\textwidth]{svgmaker_whole.png} &
    \includegraphics[width=0.45\textwidth]{svgmaker_sprtd.png} \\
    (a) Rendered output & (b) Elements separated \\
    \end{tabular}
    \caption{SVG generated by SVGMaker.io from the prompt 
    ``a penguin on ice with 16 shapes.'' The rendered output 
    (a) looks visually clean, but separating the elements 
    (b) reveals 24 paths instead of 16, several of which do 
    not correspond to recognizable visual components.}
    \label{fig:svgmaker}
\end{figure}

A systematic comparison of LLM-based SVG generation methods, 
covering both open-source and commercial systems, is a 
natural next step. Malashenko et 
al.~\cite{LeveragingLargeLanguage} provide an initial 
survey of this space. As discussed in 
Section~\ref{sec:implications}, recent work on 
chain-of-thought reasoning for SVG generation suggests 
that LLMs may develop something closer to traceable 
structural reasoning. If that happens, the analytical 
framework from this thesis could be extended to evaluate 
them on the same terms.


\chapter{Table of chosen approaches}

The following tables summarize the fourteen models analyzed in this thesis. 
Optimization-based models are listed first, followed by implicit neural methods and 
dataset-trained methods. For optimization-based methods, "trainables" refers to per-instance 
SVG parameters refined at inference time rather than learned model weights. 
Models listed as using semantic distillation (SDS or VSD) optimize vector parameters 
at inference time using gradients backpropagated from a frozen pretrained diffusion model. 
No training dataset is required for these methods. 
Dataset-driven models learn generation behavior during training and produce SVGs through 
a forward pass at inference.    

\begin{landscape}
    \thispagestyle{empty} % Optional: removes headers/footers for more space
    \centering
    \null\vfill % Vertical centering
    \small
    \section*{Inference Time Overview} % Optional title
    \begin{tabularx}{\linewidth}{l | l | l | l | l | l | X}
        \toprule
        \textbf{Model} & \textbf{Task Type} & \textbf{Inference} & \textbf{Input} & \textbf{Output} & \textbf{Supervision} & \textbf{Core Idea} \\
        \midrule
        VectorFusion & Text-to-SVG & Iterative Opt. & Text & Paths & Self-Sup & Distills diffusion priors via differentiable rasterization. \\ \hline
        SVGFusion & Text-to-SVG & Forward Dec. & Text & Paths & Supervised & Continuous latent space learning using SVGX dataset. \\ \hline
        SVGDreamer & Text-to-SVG & Iterative Opt. & Text & Paths & Self-Sup & Semantic-driven vectorization (SIVE) and Particle SDS. \\ \hline
        SVGDreamer++ & Text-to-SVG & Iterative Opt. & Text & Paths & Self-Sup & Hierarchical optimization with adaptive primitive control. \\ \hline
        T2V-NPR & Text-to-SVG & Iterative Opt. & Text & Layers/Paths & Self-Sup & Dual-branch VAE learns path latent space constraints. \\ \hline
        LayerTracer & System/Pipeline & Hybrid & Text/Img & Layers/Paths & Supervised & Models designer construction logic via DiT sequences. \\ \hline
        Im2Vec & Vectorization & Forward Dec. & Img & Paths & Self-Sup & Differentiable rasterizer allows learning from appearance. \\ \hline
        SAMVG & Vectorization & Iterative Opt. & Image & Paths & Opt.-based & Uses SAM for semantic initialization in vectorization. \\ \hline
        NIVeL & Text-to-SVG & Iterative Opt. & Text & Layers/Paths & Self-Sup & Implicit neural layers as optimization domain for SDS. \\ \hline
        NeuralSVG & Text-to-SVG & Iterative Opt. & Text & Layers/Paths & Self-Sup & Scene encoded in MLP weights for layered shapes. \\ \hline
        LIVE & Vectorization & Iterative Opt. & Image & Layers/Paths & Opt.-based & Layer-wise hierarchical optimization of vector graphs. \\ \hline
        DeepSVG & Generation & Forward Dec. & None & Paths & Supervised & Hierarchical Transformer for shape/command disentanglement. \\ \hline
        DeepIcon & Vectorization & Forward Dec. & Image & Layers/Paths & Supervised & Hierarchical network for direct image-to-SVG translation. \\ \hline
        \bottomrule
    \end{tabularx}
    \vfill\null
    \end{landscape}






\begin{landscape}
\thispagestyle{empty}
\centering
\small
\section*{Training Architecture – Optimization-Based Methods}

\begin{tabularx}{\linewidth}{
>{\raggedright\arraybackslash}p{2.3cm} |
>{\raggedright\arraybackslash}p{3cm} |
>{\raggedright\arraybackslash}p{2.2cm} |
>{\raggedright\arraybackslash}p{2cm} |
>{\raggedright\arraybackslash}p{2.5cm} |
>{\raggedright\arraybackslash}p{2.8cm} |
X}
\toprule
\textbf{Model} & \textbf{Trainables} & \textbf{Training Data} & \textbf{Init.} & \textbf{Supervision} & \textbf{Representation} & \textbf{Priors} \\
\midrule
LIVE & SVG parameters & None & Raster-driven & None & Parametric (flat) & None \\ \hline
SAMVG & None & None & Mask-based & None & Parametric (mask-derived) & SAM \\ \hline
VectorFusion & SVG parameters & None & Random / Trace & SDS (Stable Diff.) & Parametric (flat) & Stable Diff. \\ \hline
SVGDreamer & SVG parameters / LoRA & None & Random paths & VPSD & Parametric (flat) & Stable Diff. \\ \hline
SVGDream++ & SVG parameters / LoRA & None & Random + HIVE & VPSD + HIVE loss & Structured parametric & Stable Diff. + SAM \\ \hline
NeuralSVG & MLP weights + colors/ LoRA & None & Saliency-based & SDS & Implicit neural fields & Stable Diff. \\ \hline
NIVeL & MLP weights + layer colors & None & Random weights & SDS & Implicit neural fields & DeepFloyd + CLIP \\
\bottomrule
\end{tabularx}
\end{landscape}



\begin{landscape}
\thispagestyle{empty}
\centering
\small
\section*{Training Architecture – Dataset-Trained Methods}

\begin{tabularx}{\linewidth}{
>{\raggedright\arraybackslash}p{2.3cm} |
>{\raggedright\arraybackslash}p{3cm} |
>{\raggedright\arraybackslash}p{2.5cm} |
>{\raggedright\arraybackslash}p{2cm} |
>{\raggedright\arraybackslash}p{2.5cm} |
>{\raggedright\arraybackslash}p{2.8cm} |
X}
\toprule
\textbf{Model} & \textbf{Trainables} & \textbf{Training Data} & \textbf{Init.} & \textbf{Supervision} & \textbf{Representation} & \textbf{Priors} \\
\midrule
DeepSVG & VAE (Transformer) weights & SVGIcons8 & Latent sampling & Recon. + KL & Sequential path tokens & None \\ \hline
DeepIcon & Transformer weights & SVGIcons8 & CLIP embedding & Reconstruction & Sequential path tokens & CLIP encoder \\ \hline
SVGFusion & VP-VAE + VS-DiT & SVGX & Random latent & Recon. + Denoising & Continuous latent space & CLIP + DINO-v2 \\ \hline
T2V-NPR & NPR-VAE + Latents + LoRA & FIGR-8 & Random path latents & VSD & Neural Path Repr. & Stable Diff. + CLIP \\ \hline
Im2Vec & Encoder + RNN + 1D CNN & Raster images & Image encoding & Reconstruction & Parametric paths & None \\ \hline
LayerTracer & DiT + LoRA & Serpentine dataset & Diffusion noise & Denoising loss & Layered pixel sequences & Flux DiT + Florence-2 \\
\bottomrule
\end{tabularx}
\end{landscape}





\begin{landscape}
\thispagestyle{empty}
\centering
\null\vfill
\small
\section*{Reproducibility Overview}

\begin{tabularx}{\linewidth}{l | l | l | l | l | X}
\toprule
\textbf{Model} & 
\textbf{Code} & 
\textbf{Pre-trained} & 
\textbf{Hardware} & 
\textbf{Cost of time} & 
\textbf{Reproduction Feasibility} \\
\midrule

LIVE & Yes & No & RTX 3090 & 2,600--10,000 s & Feasible on high-end consumer GPU (slow). \\ \hline
SAMVG & No & No & RTX 3090 & 139--2,000 s & Re-implementation possible but requires effort. \\ \hline
VectorFusion & No & No & Not specified & Unknown & Not reproducible due to lack of public implementation. \\ \hline
SVGDreamer & Yes & No & V100 ($\geq$31GB VRAM) & 7--14 min & Requires high-VRAM GPU and LoRA pretraining. \\ \hline
SVGDreamer++ & Yes & No & A800 ($\geq$31GB VRAM) & ~7 min & Requires high-VRAM GPU and LoRA pretraining. \\ \hline
NeuralSVG & Yes & No & Not stated & Unknown & Hardware requirements unclear. \\ \hline
NIVeL & Pending & No & A100 & ~5 min & High-end GPU required. \\ \hline
T2V-NPR & Yes & VAE (unknown)& RTX 3090 & ~13 min & Feasible on high-end consumer GPU. \\ \hline
DeepSVG & Yes & VAE (unknown) & 2×1080Ti (training) & Unknown & Feasible on consumer GPU. \\ \hline
DeepIcon & No & Transformer (unknown) & Not specified & -- & Not reproducible (no official code). \\ \hline
SVGFusion & No & VP-VAE + VS-DiT (unknown) & Multi-A800 & 24--36 s & No official implementation available. \\ \hline
Im2Vec & Yes & VAE (unknown) & Not specified & Unknown & Likely feasible on consumer GPU. \\ \hline
LayerTracer & Yes & DiT + LoRA (unknown) & A100 (80GB) & 27--34 s & Enterprise GPU required. \\

\bottomrule

\end{tabularx}


\vfill\null
\end{landscape}


\section*{Architectural Pipeline Diagrams}
\label{app:pipelines}

The following diagrams visualize the generation pipelines 
of selected models analyzed in this thesis. They are 
grouped by method family to highlight where pipelines 
share components and where they diverge.

\begin{figure}[ht]
    \centering
    \includegraphics[width=\textwidth]{LIVEnSAM.png}
    \caption{Reconstruction-supervised optimization pipelines. 
    LIVE (red) and SAMVG (blue) share the same core loop: 
    Initialized SVG Params $\rightarrow$ DiffVG $\rightarrow$ 
    Rendered Image $\rightarrow$ Pixel Reconstruction Loss 
    $\rightarrow$ Target Image, with gradients flowing back 
    to update the SVG parameters. The only difference is 
    where optimization starts. LIVE initializes directly 
    from the target image through residual analysis. SAMVG 
    adds the Segment Anything Model as a preprocessing step, 
    producing semantically meaningful initialization regions 
    before the same optimization loop begins. Same loss, 
    different starting point, different structural outcome.}
    \label{fig:opt_recon_pipelines}
\end{figure}

\begin{figure}[ht]
    \centering
    \includegraphics[width=\textwidth]{OPT..png}
    \caption{Semantic-supervised optimization pipelines. 
    VectorFusion (orange), SVGDreamer (orange with blue 
    SIVE initialization), and SVGDreamer++ (purple, adding 
    hierarchical image vectorization via Grounded SAM). All 
    three share the same core optimization loop (Initialized 
    SVG Params $\rightarrow$ DiffVG $\rightarrow$ Rendered 
    Image $\rightarrow$ SDS/VPSD loss) but differ in how 
    primitives are initialized and what additional supervision 
    signals are applied. SVGDreamer adds semantic attention 
    maps for initialization. SVGDreamer++ replaces those 
    with Grounded SAM for finer part-level decomposition.}
    \label{fig:opt_semantic_pipelines}
\end{figure}

\begin{figure}[ht]
    \centering
    \includegraphics[width=\textwidth]{OPT.+NN.png}
    \caption{Neural implicit optimization pipelines. 
    NeuralSVG (blue) and NiVeL (red) both store geometry 
    in MLP weights rather than explicit B\'{e}zier 
    coordinates, but their pipelines differ. NeuralSVG 
    generates an initial layout through saliency-based 
    sampling from a diffusion model, encodes shape indices 
    via positional encoding, and optimizes MLP weights 
    using SDS loss with a LoRA adapter on Stable Diffusion. 
    NiVeL takes spatial coordinates as input, applies 
    frequency encoding, and optimizes layered implicit 
    fields using SDS with entropy loss through DeepFloyd. 
    Both produce SVG output by converting the learned 
    implicit representation into discrete vector primitives.}
    \label{fig:opt_neural_pipelines}
\end{figure}


